% !PRINTED 655
\weberpage{655}
%<<CONT
then $x_1c_1 + x_2c_2 + \cdots + x_nc_n$ is divisible by $\pi$, and, since it is a
function in $Z$, by $(z - c)$. We set:
\[
x_1c_1 + x_2c_2 + \cdots + x_nc_n = (z - c)y
\]
and obtain from (8)
\weq{9}{
\begin{aligned}
\omega &= x_1(\omega_1 - c_1) + x_2(\omega_2 - c_2) + \cdots + x_n(\omega_n - c_n) + (z - c)y\\
&= (\tau - \tau_0) + (z - c)y.
\end{aligned}
}

If now $(\tau - \tau_0)$ and $(z - c)$ had the common divisor
$\pi\delta$, then by \S\ 182, 11. one could find a function $\omega$ in $\Omega$ that is
divisible by $\pi$ and coprime to $\delta$. But equation (9) contradicts
this, and with that our theorem is proved.

Let
\weq{10}{
\Phi(t) = B_0t^m + B_1t^{m-1} + \cdots + B_{m-1}t + B_m
}
be a function of degree $m$ in $t$, and let $B_0, B_1, \ldots, B_m$ be integral
holomorphic functionals in $\overline{Z}$ with the functional variables $t_1,
t_2, \ldots, t_n$, the function satisfying the congruence condition
\weq{11}{
\Phi(\tau) \equiv 0 \quad (\mathrm{mod}\ z - c)
}
[as, for example, the function $F(t) = N(t - \tau)$ does]. Then
\[
\Phi(\tau_0) \equiv 0 \quad (\mathrm{mod}\ \pi)\,,
\]
and if we divide $\Phi(t)$ by $\tau - \tau_0$, there results
% [sic] printed $\tau - \tau_0$ (Greek tau) as the divisor of $\Phi(t)$, but division by $t - \tau_0$ is meant, as the factor $(t - \tau_0)$ in (12) shows
\weq{12}{
\Phi(t) \equiv (t - \tau_0)\Phi_1(t) \quad (\mathrm{mod}\ \pi)\,,
}
in which $\Phi_1(t)$ is a holomorphic integral functional in $\overline{Z}$ (because
$\tau_0$ is contained in $\overline{Z}$ as a \gesperrt{constant} functional). Since now in (12)
only rational functions of $z$ occur besides the functional variables,
this congruence must also be valid for the modulus $z - c$:
\weq{13}{
\Phi(t) \equiv (t - \tau_0)\Phi_1(t) \quad (\mathrm{mod}\ z - c).
}

If $\pi'$ is a second prime functional, which may also be identical
with $\pi$, and if $z - c$ is divisible by $\pi\pi'$, and $\tau \equiv \tau_0' \ (\mathrm{mod}\ \pi')$,
then $(\tau - \tau_0)\!: \pi$ is relatively prime to $\pi'$ (by 3.), and because of (11)
we have $\Phi_1(\tau_0') \equiv 0 \ (\mathrm{mod}\ \pi')$. From this one concludes
\[
\Phi_1(t) \equiv (t - \tau_0')\Phi_2(t)
\]
first for the modulus $\pi'$, but then, since the functionals
lie in $\overline{Z}$, also for the modulus $(z - c)$. We therefore have:
\weq{14}{
\Phi(t) \equiv (t - \tau_0)(t - \tau_0')\Phi_2(t) \quad (\mathrm{mod}\ z - c).
}

From this we infer the following important theorem:

4. \gesperrt{Let, according to \S\ 183, (17)}
\weq{15}{
z - c = \pi_1\pi_2\ldots\pi_n
}
%CONT>>
